\section{Conclusions and open problems}\label{sec:concl}
\subsection{What has been established}
The central question was which geometry changes along a collective orbit are exactly predictable from the arithmetic state, which depend on additional relations, and which topological consequences need cellular-attachment data. The answer assembled here is as follows.
\begin{itemize}[itemsep=2pt]
\item \emph{Exactly determined by the arithmetic state.} The realised complexes are functions of the occurrence set $D$ (\cref{thm:R3T1a,thm:A5T31}); for integers they collapse to the single integer $\sigma$ and are contractible (\cref{thm:B2T1}); occurrences are monotone along orbits (\cref{lem:M}); the next occurrence event is given by one linear congruence per candidate (\cref{thm:sieve,cor:event}) and events never stop (\cref{cor:never}); the old edges that die and the new edges that appear are read from divisor intervals (\cref{thm:A5T32}).
\item \emph{Depending on additional relations.} Equal occurrence sets can carry different primitive relations (\cref{thm:R3T1,thm:R3T2}); relations can be dormant and become active without being created (\cref{thm:A4T1,thm:A4T2}); whether an activated target lies on the orbit is a reachability question with a proved obstruction and exact finite examples (\cref{thm:A5T11,thm:A4T3c,thm:A4T3c375}).
\item \emph{Needing cellular-attachment data.} The change of Betti numbers follows the zig-zag identity, with local terms and global ranks (\cref{thm:zigzag}); a change always has a cause in dead or new cells (\cref{thm:cause}); at graph level the update is exact (\cref{thm:A5T44}); activation does not force a topological change and does not always raise homology (\cref{prop:A4T4c,thm:A5T7}).
\item \emph{Towers.} Both towers are nested with no dead cell (\cref{thm:A5T51}); an eventual quasi-polynomial law exists with period dividing the order of the growing prime, with uncontrolled threshold (\cref{thm:A5T53}); the exact laws are finite-exact to $k\le140$ and are not promoted (\cref{thm:A5T6}).
\item \emph{Correspondences.} Between the collective and the generative systems the exact statements are transport formulas, the snapshot relation $C^3=4x=T_1(x)$, the one-step and scaling obstructions, a coincidence family and the classification of formal multipliers on the two towers; no universal correspondence is established (\cref{sec:corr}).
\end{itemize}
Corrections carried out in this paper are the following: in B2-T7 only $G_\infty=1\Rightarrow$ rigidity holds, the $G=5$ tower being a counterexample to the converse (\cref{prop:B2T7}); in B2-T9, $\sigma$ strictly increases and $D$ changes on every edge, but inclusion $D(m)\subseteq D(m')$ fails (\cref{prop:B2T9}); the inherited proof of the cause criterion was repaired (\cref{thm:cause}); the loss step of B2-T4 is the even one (\cref{thm:B2T4}); the entry-time statistic for the regime-I descent is $12$ in the range checked (\cref{rem:entry}).

\subsection{Limits of this paper}
All census and tower results are finite-exact and are stated with their ranges; no extrapolated curve is plotted. Inherited verifier suites and inherited censuses (the $993$-pair reachability table, the $474$-event table, the census of $1{,}920$ orbits, the $20{,}000$-candidate rigid sample, the $16{,}028$-pair birth-endpoint scan) are \stat{cited} and not re-run. The ownership of results shared with Paper~II is recorded in Appendix~\ref{app:crosswalk}. The programme sources~\cite{FoundI,GEOI} are cited at their pinned commits and were not re-fetched.

\subsection{Open problems}
\begin{enumerate}[label=\textbf{P\arabic*.},itemsep=2pt]
\item \emph{Relations beyond $D$} (\cref{prob:R3T2p}): classify the minimal descriptor $\Phi(D,J)$ for boxes $q\ne\mathbf 1^k$ and three or more primes.
\item \emph{Reachability} (\cref{prob:A4T3f}): are infinitely many members of the activation family on the orbit of their source? Is every non-blocked member reached? A closed form for the landing step?
\item \emph{Type predicate}: no arithmetic predicate of the state is known that predicts the type (V, II, III, IV) of an occurrence event; only the form of the answer, determined by $(K(D),N)$, is proved (\cref{thm:A5T7}).
\item \emph{Dead edges and $\beta_1$}: in the census no event with a dead edge raises $\beta_1$; in the abstract calculus it can (\cref{cx:abstract}); can it along a collective orbit?
\item \emph{Higher-dimensional predictor}: a rebuild-free prediction of the Betti vector when $2$-cells and higher cells are present (\cref{thm:A5T44} is graph-level only).
\item \emph{Threshold and mechanism for the towers}: an effective bound for $k_1$ in \cref{thm:A5T53}; an explanation of why $G=11$ has growing $\beta_3$ and $G=5$ none (\cref{prop:A4T56}).
\item \emph{Periodicity}: is the stage sequence periodic for every persistent orbit? For $\ell=29$ the periodic pair $(P,\rho^2)=(35,900)$ is verified, not proved (\cref{ex:l29}); the $(3,4),(8,9),(35,900)$ census is not a theorem.
\item \emph{Mixed-index forcing for $k>2$} primes (B2-T5), \emph{completeness of single-edge coincidences} (B2-T12), and the connection between history complexes and orbits beyond the obstruction (\cref{prop:A4T7}).
\end{enumerate}

\subsection{Place in the programme}
Papers~I--III together prepare, but do not contain, the synthesis of Manuscript~IV, \emph{The Geometry of Generation in the UNNS Substrate: A Mathematical Synthesis}. Nothing in this paper claims a universal organising law; the results above are statements about the collective map, the selected towers and the restricted correspondences, each with the status and range recorded in the theorem map (Appendix~\ref{app:crosswalk}). Transport and dimensional obstructions are the subject of Paper~I~\cite{PaperI}, history complexes that of Paper~II~\cite{PaperII}, and the geometric framework on which all three rely is that of~\cite{GEOI}.
